\documentclass[12pt]{article}
\usepackage[utf8]{inputenc}
\usepackage[T1]{fontenc}
\usepackage{hyperref}
\title{Limi Language Resources for K-12 Curriculum Development}
\author{Compiled by Claude Sonnet 5.5 (claude-sonnet-5-5) for Andr\'as Kornai}
\date{}
\begin{document}
\maketitle

\section{Introduction}
Limi (ISO 639-3: ylm) is a Loloish Sino-Tibetan language of western Yunnan, China. It is spoken in Yongde, Fengqing and Yun counties. Speaker estimates range from about 20,000 (Yang 2017) to 29,000 (2002) and 38,000 (Joshua Project). Speakers are officially counted within the Yi nationality. Asia Harvest describes the language as well preserved and in active use. Written resources appear to be very limited.

\section{Primary Educational Resources}
No Limi primers, textbooks or school programmes were found in search results.

\section{Linguistic and Reference Materials}
\begin{itemize}
\item Wikipedia, ``Limi language''. It says Limi is probably closest to Lolopo and has many Lalo loanwords. \url{https://en.wikipedia.org/wiki/Limi_language}
\item Glottolog entry for Limi. \url{https://glottolog.org/resource/languoid/id/limi1243.bigmap.html}
\end{itemize}

\section{Religious and Community Resources}
\begin{itemize}
\item Joshua Project, Limi in China. It lists Bible translation as needed. \url{https://joshuaproject.net/people_groups/18554}
\item Joshua Project language page for ylm. \url{https://joshuaproject.net/languages/ylm}
\item Asia Harvest people-group profile of the Limi (PDF). \url{https://people-groups.asiaharvest.org/people-group-profiles/yunnan/limi.pdf}
\end{itemize}

\section{Digital and Audio Resources}
No audio, video or software resources were found.

\section{Curriculum Development Context}
The language is in active community use, which favors native-medium work. Because no written standard or school materials were found, orthography development and materials creation would start from scratch, with community involvement.

\end{document}